\documentclass[a4paper,12pt,twoside]{memoir}

% Castellano
\usepackage[spanish,es-tabla]{babel}
\selectlanguage{spanish}
\usepackage[utf8]{inputenc}
\usepackage[T1]{fontenc}
\usepackage{lmodern} % scalable font
\usepackage{microtype}
\usepackage{placeins}
\usepackage{longtable}

\RequirePackage{booktabs}
\RequirePackage[table]{xcolor}
\RequirePackage{xtab}
\RequirePackage{multirow}


% Links
\PassOptionsToPackage{hyphens}{url}\usepackage[colorlinks]{hyperref}
\hypersetup{
	allcolors = {red}
}

% Ecuaciones
\usepackage{amsmath}

% Rutas de fichero / paquete
\newcommand{\ruta}[1]{{\sffamily #1}}

% Párrafos
\nonzeroparskip

% Huérfanas y viudas
\widowpenalty100000
\clubpenalty100000

% Evitar solapes en el header
\nouppercaseheads

% Imagenes
\usepackage{graphicx}
\newcommand{\imagen}[2]{
	\begin{figure}[!h]
		\centering
		\includegraphics[width=0.9\textwidth]{#1}
		\caption{#2}\label{fig:#1}
	\end{figure}
	\FloatBarrier
}

\newcommand{\imagenflotante}[2]{
	\begin{figure}%[!h]
		\centering
		\includegraphics[width=0.9\textwidth]{#1}
		\caption{#2}\label{fig:#1}
	\end{figure}
}

\newcommand{\imagenTam}[3]{
	\begin{figure}[!h]
		\centering
		\includegraphics[width=#3\textwidth]{#1}
		\caption{#2}\label{fig:#1}
	\end{figure}
	\FloatBarrier
}


% El comando \figura nos permite insertar figuras comodamente, y utilizando
% siempre el mismo formato. Los parametros son:
% 1 -> Porcentaje del ancho de página que ocupará la figura (de 0 a 1)
% 2 --> Fichero de la imagen
% 3 --> Texto a pie de imagen
% 4 --> Etiqueta (label) para referencias
% 5 --> Opciones que queramos pasarle al \includegraphics
% 6 --> Opciones de posicionamiento a pasarle a \begin{figure}
\newcommand{\figuraConPosicion}[6]{%
  \setlength{\anchoFloat}{#1\textwidth}%
  \addtolength{\anchoFloat}{-4\fboxsep}%
  \setlength{\anchoFigura}{\anchoFloat}%
  \begin{figure}[#6]
    \begin{center}%
      \Ovalbox{%
        \begin{minipage}{\anchoFloat}%
          \begin{center}%
            \includegraphics[width=\anchoFigura,#5]{#2}%
            \caption{#3}%
            \label{#4}%
          \end{center}%
        \end{minipage}
      }%
    \end{center}%
  \end{figure}%
}

%
% Comando para incluir imágenes en formato apaisado (sin marco).
\newcommand{\figuraApaisadaSinMarco}[5]{%
  \begin{figure}%
    \begin{center}%
    \includegraphics[angle=90,height=#1\textheight,#5]{#2}%
    \caption{#3}%
    \label{#4}%
    \end{center}%
  \end{figure}%
}
% Para las tablas
\newcommand{\otoprule}{\midrule [\heavyrulewidth]}
%
% Nuevo comando para tablas pequeñas (menos de una página).
\newcommand{\tablaSmall}[5]{%
 \begin{table}
  \begin{center}
   \rowcolors {2}{gray!35}{}
   \begin{tabular}{#2}
    \toprule
    #4
    \otoprule
    #5
    \bottomrule
   \end{tabular}
   \caption{#1}
   \label{tabla:#3}
  \end{center}
 \end{table}
}

%
%Para el float H de tablaSmallSinColores
\usepackage{float}

%
% Nuevo comando para tablas pequeñas (menos de una página).
\newcommand{\tablaSmallSinColores}[5]{%
 \begin{table}[H]
  \begin{center}
   \begin{tabular}{#2}
    \toprule
    #4
    \otoprule
    #5
    \bottomrule
   \end{tabular}
   \caption{#1}
   \label{tabla:#3}
  \end{center}
 \end{table}
}

\newcommand{\tablaApaisadaSmall}[5]{%
\begin{landscape}
  \begin{table}
   \begin{center}
    \rowcolors {2}{gray!35}{}
    \begin{tabular}{#2}
     \toprule
     #4
     \otoprule
     #5
     \bottomrule
    \end{tabular}
    \caption{#1}
    \label{tabla:#3}
   \end{center}
  \end{table}
\end{landscape}
}

%
% Nuevo comando para tablas grandes con cabecera y filas alternas coloreadas en gris.
\newcommand{\tabla}[6]{%
  \begin{center}
    \tablefirsthead{
      \toprule
      #5
      \otoprule
    }
    \tablehead{
      \multicolumn{#3}{l}{\small\sl continúa desde la página anterior}\\
      \toprule
      #5
      \otoprule
    }
    \tabletail{
      \hline
      \multicolumn{#3}{r}{\small\sl continúa en la página siguiente}\\
    }
    \tablelasttail{
      \hline
    }
    \bottomcaption{#1}
    \rowcolors {2}{gray!35}{}
    \begin{xtabular}{#2}
      #6
      \bottomrule
    \end{xtabular}
    \label{tabla:#4}
  \end{center}
}

%
% Nuevo comando para tablas grandes con cabecera.
\newcommand{\tablaSinColores}[6]{%
  \begin{center}
    \tablefirsthead{
      \toprule
      #5
      \otoprule
    }
    \tablehead{
      \multicolumn{#3}{l}{\small\sl continúa desde la página anterior}\\
      \toprule
      #5
      \otoprule
    }
    \tabletail{
      \hline
      \multicolumn{#3}{r}{\small\sl continúa en la página siguiente}\\
    }
    \tablelasttail{
      \hline
    }
    \bottomcaption{#1}
    \begin{xtabular}{#2}
      #6
      \bottomrule
    \end{xtabular}
    \label{tabla:#4}
  \end{center}
}

%
% Nuevo comando para tablas grandes sin cabecera.
\newcommand{\tablaSinCabecera}[5]{%
  \begin{center}
    \tablefirsthead{
      \toprule
    }
    \tablehead{
      \multicolumn{#3}{l}{\small\sl continúa desde la página anterior}\\
      \hline
    }
    \tabletail{
      \hline
      \multicolumn{#3}{r}{\small\sl continúa en la página siguiente}\\
    }
    \tablelasttail{
      \hline
    }
    \bottomcaption{#1}
  \begin{xtabular}{#2}
    #5
   \bottomrule
  \end{xtabular}
  \label{tabla:#4}
  \end{center}
}



\definecolor{cgoLight}{HTML}{EEEEEE}
\definecolor{cgoExtralight}{HTML}{FFFFFF}

%
% Nuevo comando para tablas grandes sin cabecera.
\newcommand{\tablaSinCabeceraConBandas}[5]{%
  \begin{center}
    \tablefirsthead{
      \toprule
    }
    \tablehead{
      \multicolumn{#3}{l}{\small\sl continúa desde la página anterior}\\
      \hline
    }
    \tabletail{
      \hline
      \multicolumn{#3}{r}{\small\sl continúa en la página siguiente}\\
    }
    \tablelasttail{
      \hline
    }
    \bottomcaption{#1}
    \rowcolors[]{1}{cgoExtralight}{cgoLight}

  \begin{xtabular}{#2}
    #5
   \bottomrule
  \end{xtabular}
  \label{tabla:#4}
  \end{center}
}




\graphicspath{ {./img/} }

% Capítulos
\chapterstyle{bianchi}
\newcommand{\capitulo}[2]{
	\setcounter{chapter}{#1}
	\setcounter{section}{0}
	\setcounter{figure}{0}
	\setcounter{table}{0}
	\chapter*{#2}
	\addcontentsline{toc}{chapter}{#2}
	\markboth{#2}{#2}
}

% Apéndices
\renewcommand{\appendixname}{Apéndice}
\renewcommand*\cftappendixname{\appendixname}

\newcommand{\apendice}[1]{
	%\renewcommand{\thechapter}{A}
	\chapter{#1}
}

\renewcommand*\cftappendixname{\appendixname\ }

% Formato de portada
\makeatletter
\usepackage{xcolor}
\newcommand{\tutor}[1]{\def\@tutor{#1}}
\newcommand{\course}[1]{\def\@course{#1}}
\definecolor{cpardoBox}{HTML}{E6E6FF}
\def\maketitle{
  \null
  \thispagestyle{empty}
  % Cabecera ----------------
\noindent\includegraphics[width=\textwidth]{cabecera}\vspace{1cm}%
  \vfill
  % Título proyecto y escudo informática ----------------
  \colorbox{cpardoBox}{%
    \begin{minipage}{.8\textwidth}
      \vspace{.5cm}\Large
      \begin{center}
      \textbf{TFG del Grado en Ingeniería Informática}\vspace{.6cm}\\
      \textbf{\LARGE\@title{}}
      \end{center}
      \vspace{.2cm}
    \end{minipage}

  }%
  \hfill\begin{minipage}{.20\textwidth}
    \includegraphics[width=\textwidth]{escudoInfor}
  \end{minipage}
  \vfill
  % Datos de alumno, curso y tutores ------------------
  \begin{center}%
  {%
    \noindent\LARGE
    Presentado por \@author{}\\ 
    en Universidad de Burgos --- \@date{}\\
    Tutores: \@tutor{}\\
  }%
  \end{center}%
  \null
  \cleardoublepage
  }
\makeatother


% Datos de portada
\title{Aplicación Web para Gestión de Datos Genómicos de Cánceres de Pulmón}
\author{Carlos Ortúñez Rojo}
\tutor{Dr. Jesús Manuel Maudes Raedo y Dr. José Francisco Díez Pastor}
\date{\today}

\begin{document}

\maketitle



\cleardoublepage



%%%%%%%%%%%%%%%%%%%%%%%%%%%%%%%%%%%%%%%%%%%%%%%%%%%%%%%%%%%%%%%%%%%%%%%%%%%%%%%%%%%%%%%%



\frontmatter


\clearpage

% Indices
\tableofcontents

\clearpage

\listoffigures

\clearpage

\listoftables

\clearpage

\mainmatter

\appendix

\apendice{Plan de Proyecto Software}

\section{Introducción}

En este apartado se va a tratar la planificación temporal que se ha llevado a cabo para el proyecto, así como un estudio de viabilidad tanto económica como legal.

\section{Planificación temporal}



\section{Estudio de viabilidad}

\subsection{Viabilidad económica}

\subsection{Viabilidad legal}



\apendice{Especificación de Requisitos}

\section{Introducción}

En este anexo se van a especificar los requisitos funcionales, los requisitos no funcionales y los casos de uso. Estos requisitos no han sido establecidos de forma trivial sino que se han obtenido de las reuniones iniciales con los oncólogos.

\section{Objetivos generales}

Los objetivos generales acordados con los oncólogos en las primeras reuniones son los siguientes:

\begin{itemize}
    \item Poder gestionar los usuarios que tiene acceso a la web.
    \item Poder gestionar los datos de los pacientes.
	\item Tener una forma de visualizar estos datos de manera sencilla.
	\item La web ha de ser sencilla e intuitiva de usar.
\end{itemize}

\section{Catalogo de requisitos}

En esta sección se van a enumerar tanto los requisitos funcionales como los no funcionales, haciendo una breve explicación de ellos.

\subsection{Requisitos funcionales}

\begin{itemize}
    \item \textbf{RF-1 Sistema de autenticación de los usuario:} Solo los usuarios autorizados podrán acceder al sistema
    \begin{itemize}
        \item \textbf{RF-1.1 Identificación de usuario:} Los usuarios se identificaran con un correo y una contraseña únicos.
        \item \textbf{RF-1.2 Uso de dos roles para los usuarios:} Existirá el rol administrador y el rol oncólogo.
    	\item \textbf{RF-1.3 Gestión de usuarios solo para administradores:} Solo los usuarios con el rol administrador tendrán acceso a la gestión de usuarios.
    \end{itemize}
    \item \textbf{RF-2 Sistema de gestión de usuarios:} El administrador podrá gestionar los usuarios pudiendo añadir, modificar y eliminar usuarios.
    \begin{itemize}
        \item \textbf{RF-2.1 Añadir usuario:} El administrador podrá añadir nuevos usuarios al sistema.
        \item \textbf{RF-2.2 Modificar datos usuario:} El administrador podrá modificar los datos de los usuarios.
        \begin{itemize}
            \item \textbf{RF-2.2.1 Modificar datos personales usuario:} El administrador podrá modificar los datos personales de los usuarios.
            \item \textbf{RF-2.2.2 Modificar rol usuario:} El administrador podra modificar el rol de los usuarios.
        \end{itemize}
        \item \textbf{RF-2.3 Eliminar usuario:} El administrador podrá eliminar usuarios que no tengan el rol de administrador.
    \end{itemize}
    \item \textbf{RF-3 Sistema de gestión de datos personales:} Los usuarios podrán
    modificar sus datos personales.
    \begin{itemize}
        \item \textbf{RF-3.1 Modificar datos personales:} El usuario podrá modificar sus datos personales con la única restricción que el correo no debe de estar ya registrado en la web.
        \item \textbf{RF-3.2 Modificar contraseña:} El usuario podrá modificar su contraseña siempre y cuando conozca su contraseña actual.
    \end{itemize}
    \item \textbf{RF-4 Sistema de gestión de pacientes:} Los usuarios podrán
    gestionar los pacientes pudiendo añadir, modificar y eliminar pacientes.
    \begin{itemize}
        \item \textbf{RF-4.1 Añadir pacientes:} Los usuarios podrán añadir nuevos pacientes a la web.
        \item \textbf{RF-4.2 Añadir datos paciente:} Los usuarios podrán añadir datos de los pacientes.
        \begin{itemize}
            \item \textbf{RF-4.2.1 Añadir enfermedad:} El usuario podrá añadir una enfermedad por cada paciente
            \item \textbf{RF-4.2.2 Añadir antecedentes:} El usuario podrá añadir antecedentes a un paciente.
            \item \textbf{RF-4.2.3 Añadir tratamientos:} El usuario podrá añadir tratamientos a un paciente.
            \item \textbf{RF-4.2.4 Añadir reevaluaciones:} El usuario podrá añadir reevaluaciones a un paciente.
            \item \textbf{RF-4.2.5 Añadir seguimientos:} El usuario podrá añadir seguimientos a un paciente.
            \item \textbf{RF-4.2.6 Añadir comentarios:} El usuario podrá añadir comentarios a un paciente.
        \end{itemize}
        \item \textbf{RF-4.3 Modificar paciente:} Los usuarios podrán modificar los datos de los pacientes.
        \begin{itemize}
            \item \textbf{RF-4.3.1 Modificar datos demográficos:} El usuario podrá modificar los datos demográficos de un paciente.
            \item \textbf{RF-4.3.2 Modificar datos enfermedad:} El usuario podrá modificar los datos de la enfermedad de un paciente.
            \item \textbf{RF-4.3.3 Modificar datos antecedentes:} El usuario podrá modificar los datos de los antecedentes del paciente..
            \item \textbf{RF-4.3.4 Modificar datos tratamientos:} El usuario podrá modificar los datos de los tratamientos del paciente..
            \item \textbf{RF-4.3.5 Modificar datos reevaluaciones:} El usuario podrá modificar los datos de las reevaluaciones del paciente.
            \item \textbf{RF-4.3.6 Modificar datos seguimientos:} El usuario podrá modificar los datos de los seguimientos del paciente.
            \item \textbf{RF-4.3.7 Modificar comentarios:} El usuario podrá modificar los datos demográficos de los comentarios realizados al paciente.
        \end{itemize}
        \item \textbf{RF-4.4 Eliminar paciente:} Los usuarios podrán eliminar pacientes y todos sus datos correspondientes.
    \end{itemize}
	\item \textbf{RF-5 Sistema de visualización de datos:} Los usuarios podrán ver los datos de los pacientes de diferentes maneras.
	\begin{itemize}
	    \item \textbf{RF-5.1 Visualización de datos generales:} Los usuarios podrán visualizar los datos generales permitiendo la división mediante 2 atributos.
	    \item \textbf{RF-5.2 Visualización de datos numéricos:} Los usuarios podrán visualizar los datos numéricos y distinta información estadística.
	    \item \textbf{RF-5.3 Visualización de datos individualmente:}  Los usuarios podrán visualizar los datos individuales de los pacientes mediante tablas.
	    \item \textbf{RF-5.4 Exportación de los datos:} Los usuarios podrán exportar los datos de todos los pacientes a un fichero Excel.
	    \item \textbf{RF-5.5 Visualización de datos en la gráfica kaplan meier:} Los usuarios podrán visualizar los datos de supervivencia de los pacientes en una gráfica kaplan meier.
	\end{itemize}
	\item \textbf{RF-6 Creación de base de datos sintética:} Los usuarios podrán crear una base de datos sintética seleccionando distintos parámetros para los atributos numéricos.
    \item \textbf{RF-7 Filtrado de pacientes:} Los usuarios podrán filtrar pacientes según sus atributos.
	\begin{itemize}
	    \item \textbf{RF-7.1 Filtrado de pacientes tabla:} Los usuarios podrán filtrar los pacientes por atributos y visualizarlos en la tabla general de pacientes.
	    \item \textbf{RF-7.2 Filtrado de pacientes y visualización de datos generales:} Los usuarios podrán filtrar los pacientes por atributos y visualizarlos sus estadísticas en la gráfica y en la tabla.
	    \item \textbf{RF-7.3 Filtrado de pacientes y visualización de datos numéricos:} Los usuarios podrán filtrar los pacientes por atributos y visualizarlos su información estadística.
	    \item \textbf{RF-7.4 Filtrado de pacientes y visualización de gráfica kaplan meier:} Los usuarios podrán filtrar los pacientes por atributos y visualizarlos su gráfica de supervivencia.
	\end{itemize}
\end{itemize}

\subsection{Requisitos no funcionales}

\begin{itemize}
    \item \textbf{RNF-1 Usabilidad:} La aplicación ha de ser intuitiva y fácil de usar para el usuario.
    \item \textbf{RNF-2 Seguridad:} Los datos de los pacientes solo podrán ser observador por personas con acceso a la aplicación.
    \item \textbf{RNF-3 Confidencialidad:} Los datos sensibles de los pacientes han de estar encriptados en la base de datos.
    \item \textbf{RNF-4 Rendimiento:} Los tiempos de ejecución y los tiempos de carga han de ser aceptables.
    \item \textbf{RNF-5 Mantenibilidad:} La aplicación ha de ser de fácil mantenimiento y actualización.
    \item \textbf{RNF-6 Compatibilidad:} La aplicación ha de poder ser usada por los principales navegadores: Edge, Chrome, Firefox, etc.
    \item \textbf{RNF-7 Facilidad de despliegue:} La aplicación ha de ser fácil de desplegar en un servidor o en local.
\end{itemize}

\section{Especificación de requisitos}

En esta sección se van a desarrollar los casos de uso relacionados con los requisitos funcionales comentados anteriormente, con sus correspondientes diagramas.

\subsection{Actores}

En los casos de uso se pueden distinguir entre dos actores:

\begin{itemize}
    \item \textbf{Administrador:} Usuario con acceso a la gestión de usuarios, gestión de pacientes y visualización.
    \item \textbf{Oncólogo:} Usuario con acceso a la gestión de pacientes y visualización.
\end{itemize}

\subsection{Diagramas de casos de uso}

\subsubsection{Nivel 0}

\imagen{Nivel0}{Diagrama de casos de uso nivel 0.}

\newpage

\subsubsection{Nivel 1}

\imagen{gestion_usuarios}{Diagrama de casos de uso nivel 1, gestión de usuarios.}

\imagen{gestionar_pacientes}{Diagrama de casos de uso nivel 1, gestión de usuarios.}

\imagen{gestionar_datos}{Diagrama de casos de uso nivel 1, gestión de datos personales.}

\imagen{visualizar_datos}{Diagrama de casos de uso nivel 1, visualizar datos.}

\imagen{visualizar_datos}{Diagrama de casos de uso nivel 1, filtrar datos.}

\subsection{Especificación de casos de uso}

\tablaSmallSinColores{CU-1: Iniciar sesión.}{p{3cm} p{.75cm} p{9cm}}{tablaCU1}{
	\multicolumn{3}{p{10.25cm}}{CU-1: Iniciar sesión}\\}
{
	 	 Descripción &\multicolumn{2}{p{10.25cm}}{El usuario se identifica en el sistema mediante el correo y la contraseña} \\\cline{2-3} 
	 	 Precondiciones & \multicolumn{2}{p{10.25cm}}{No existe una sesión activa en ese ordenador}\\\cline{2-3} 
	 	 Requisitos asociados & \multicolumn{2}{p{10.25cm}}{RF-1, RF-1.1}\\\cline{2-3} 
	 	 Usuario & \multicolumn{2}{p{10.25cm}}{Anónimo}\\\cline{2-3} 
	 	 \multirow{3}{3.5cm}{Acciones}  & Paso & Acción \\
	 	 & 1. & El usuario accede a la url del login\\
	 	 & 2. & El envía sus datos al servidor mediante el formulario\\
	 	 & 3. & El servidor valida los datos y da acceso a la web \\\cline{2-3} 
	 	 \multirow{2}{3.5cm}{Excepciones} & Id & Excepción \\
	 	 & 1. & Que el usuario o la contraseña sean incorrectos\\
	 	 & 2. & Campo en blanco\\\cline{2-3}
	 	 Importancia & \multicolumn{2}{p{10.25cm}}{Alta}\\
}

\tablaSmallSinColores{CU-2: Gestionar usuarios.}{p{3cm} p{.75cm} p{9cm}}{tablaCU2}{
	\multicolumn{3}{p{10.25cm}}{CU-2: Gestionar usuarios}\\}
{
	 	 Descripción &\multicolumn{2}{p{10.25cm}}{El administrador podrá modificar, eliminar y crear usuarios} \\\cline{2-3} 
	 	 Precondiciones & \multicolumn{2}{p{10.25cm}}{Administrador con sesión iniciada}\\\cline{2-3} 
	 	 Requisitos asociados & \multicolumn{2}{p{10.25cm}}{RF-1.3, RF-2}\\\cline{2-3} 
	 	 Usuario & \multicolumn{2}{p{10.25cm}}{Administrador}\\\cline{2-3} 
	 	 \multirow{3}{3.5cm}{Acciones}  & Paso & Acción \\
	 	 & 1. & El administrador accede en el menú a la parte de gestión de usuarios\\\cline{2-3} 
	 	 \multirow{2}{3.5cm}{Excepciones} & Id & Excepción \\
	 	 & 1. & Ninguna\\\cline{2-3}
	 	 Importancia & \multicolumn{2}{p{10.25cm}}{Alta}\\
}

\tablaSmallSinColores{CU-2.1: Añadir usuarios.}{p{3cm} p{.75cm} p{9cm}}{tablaCU2.1}{
	\multicolumn{3}{p{10.25cm}}{CU-2.1: Añadir usuarios}\\}
{
	 	 Descripción &\multicolumn{2}{p{10.25cm}}{El administrador podrá añadir un nuevo usuario} \\\cline{2-3} 
	 	 Precondiciones & \multicolumn{2}{p{10.25cm}}{Administrador con sesión iniciada}\\\cline{2-3} 
	 	 Requisitos asociados & \multicolumn{2}{p{10.25cm}}{RF-2.1}\\\cline{2-3} 
	 	 Usuario & \multicolumn{2}{p{10.25cm}}{Administrador}\\\cline{2-3} 
	 	 \multirow{3}{3.5cm}{Acciones}  & Paso & Acción \\
	 	 & 1. & El administrador selecciona añadir nuevo usuario\\
	 	 & 2. & Se rellenan los campos con los datos del usuario \\
	 	 & 3. & Se asigna un rol al usuario\\
	 	 & 4. & Se confirma el formulario\\\cline{2-3} 
	 	 \multirow{2}{3.5cm}{Excepciones} & Id & Excepción \\
	 	 & 1. & Que el correo ya exista en la base de datos\\
	 	 & 2. & Que las contraseñas no coincidan\\
	 	 & 3. & Campo en blanco\\\cline{2-3}
	 	 Importancia & \multicolumn{2}{p{10.25cm}}{Alta}\\
}

\tablaSmallSinColores{CU-2.2: Modificar usuarios.}{p{3cm} p{.75cm} p{9cm}}{tablaCU2.2}{
	\multicolumn{3}{p{10.25cm}}{CU-2.2: Modificar usuarios}\\}
{
	 	 Descripción &\multicolumn{2}{p{10.25cm}}{El administrador podrá modificar los datos de un usuario existente} \\\cline{2-3} 
	 	 Precondiciones & \multicolumn{2}{p{10.25cm}}{Administrador con sesión iniciada}\\\cline{2-3} 
	 	 Requisitos asociados & \multicolumn{2}{p{10.25cm}}{RF-2.2, RF-2.2.1, RF-2.2.2}\\\cline{2-3} 
	 	 Usuario & \multicolumn{2}{p{10.25cm}}{Administrador}\\\cline{2-3} 
	 	 \multirow{3}{3.5cm}{Acciones}  & Paso & Acción \\
	 	 & 1. & El administrador selecciona modificar en el usuario cuyos datos se quieran\\
	 	 & 2. & Se cambian los campos que se quieran modificar \\
	 	 & 3. & Se confirma el formulario\\\cline{2-3} 
	 	 \multirow{2}{3.5cm}{Excepciones} & Id & Excepción \\
	 	 & 1. & Que el correo ya exista en la base de datos\\
	 	 & 2. & Campo en blanco\\\cline{2-3}
	 	 Importancia & \multicolumn{2}{p{10.25cm}}{Media}\\
}

\tablaSmallSinColores{CU-2.3: Eliminar usuarios.}{p{3cm} p{.75cm} p{9cm}}{tablaCU2.3}{
	\multicolumn{3}{p{10.25cm}}{CU-2.3: Eliminar usuarios}\\}
{
	 	 Descripción &\multicolumn{2}{p{10.25cm}}{El administrador podrá eliminar un usuario existente} \\\cline{2-3} 
	 	 Precondiciones & \multicolumn{2}{p{10.25cm}}{Administrador con sesión iniciada}\\\cline{2-3} 
	 	 Requisitos asociados & \multicolumn{2}{p{10.25cm}}{RF-2.3}\\\cline{2-3} 
	 	 Usuario & \multicolumn{2}{p{10.25cm}}{Administrador}\\\cline{2-3} 
	 	 \multirow{3}{3.5cm}{Acciones}  & Paso & Acción \\
	 	 & 1. & El administrador selecciona eliminar en el usuario que se quiere borrar\\
	 	 & 2. & Se confirma la eliminación \\\cline{2-3} 
	 	 \multirow{2}{3.5cm}{Excepciones} & Id & Excepción \\
	 	 & 1. & Ninguna\\\cline{2-3} 
	 	 Importancia & \multicolumn{2}{p{10.25cm}}{Alta}\\
}

\tablaSmallSinColores{CU-3: Gestionar datos personales.}{p{3cm} p{.75cm} p{9cm}}{tablaCU3}{
	\multicolumn{3}{p{10.25cm}}{CU-3: Gestionar datos personales}\\}
{
	 	 Descripción &\multicolumn{2}{p{10.25cm}}{Los usuarios podrán modificar sus datos personales} \\\cline{2-3} 
	 	 Precondiciones & \multicolumn{2}{p{10.25cm}}{Usuario con sesión iniciada}\\\cline{2-3} 
	 	 Requisitos asociados & \multicolumn{2}{p{10.25cm}}{RF-3, RF-3.1}\\\cline{2-3} 
	 	 Usuario & \multicolumn{2}{p{10.25cm}}{Administrador y oncólogo}\\\cline{2-3} 
	 	 \multirow{3}{3.5cm}{Acciones}  & Paso & Acción \\
	 	 & 1. & El usuario selecciona en el menú modificar datos personales\\
	 	 & 2. & Se cambian los campos que se quieren modificar \\\cline{2-3} 
	 	 \multirow{2}{3.5cm}{Excepciones} & Id & Excepción \\
	 	 & 1. & Que el correo ya exista en la base de datos\\
	 	 & 2. & Campo en blanco\\\cline{2-3} 
	 	 Importancia & \multicolumn{2}{p{10.25cm}}{Media}\\
}

\tablaSmallSinColores{CU-3.1: Modificar contraseña.}{p{3cm} p{.75cm} p{9cm}}{tablaCU3.1}{
	\multicolumn{3}{p{10.25cm}}{CU-3.1: Modificar contraseña}\\}
{
	 	 Descripción &\multicolumn{2}{p{10.25cm}}{Los usuarios podrán modificar su contraseña} \\\cline{2-3} 
	 	 Precondiciones & \multicolumn{2}{p{10.25cm}}{Usuario con sesión iniciada}\\\cline{2-3} 
	 	 Requisitos asociados & \multicolumn{2}{p{10.25cm}}{RF-3.2}\\\cline{2-3} 
	 	 Usuario & \multicolumn{2}{p{10.25cm}}{Administrador y oncólogo}\\\cline{2-3} 
	 	 \multirow{3}{3.5cm}{Acciones}  & Paso & Acción \\
	 	 & 1. & El usuario selecciona en el menú modificar contraseña\\
	 	 & 2. & Se introduce la contraseña antigua y la nueva contraseña \\
	 	 & 3. & Se confirma el formulario\\\cline{2-3} 
	 	 \multirow{2}{3.5cm}{Excepciones} & Id & Excepción \\
	 	 & 1. & Que las contraseñas nuevas no coincidan\\
	 	 & 2. & Que la contraseña antigua no coincida\\
	 	 & 3. & Campo en blanco.\\\cline{2-3} 
	 	 Importancia & \multicolumn{2}{p{10.25cm}}{Alta}\\
}

\tablaSmallSinColores{CU-4: Gestionar pacientes.}{p{3cm} p{.75cm} p{9cm}}{tablaCU4}{
	\multicolumn{3}{p{10.25cm}}{CU-4: Gestionar pacientes}\\}
{
	 	 Descripción &\multicolumn{2}{p{10.25cm}}{Los usuarios podrán añadir, modificar y eliminar pacientes} \\\cline{2-3} 
	 	 Precondiciones & \multicolumn{2}{p{10.25cm}}{Usuario con sesión iniciada}\\\cline{2-3} 
	 	 Requisitos asociados & \multicolumn{2}{p{10.25cm}}{RF-4}\\\cline{2-3} 
	 	 Usuario & \multicolumn{2}{p{10.25cm}}{Administrador y oncólogo}\\\cline{2-3} 
	 	 \multirow{3}{3.5cm}{Acciones}  & Paso & Acción \\
	 	 & 1. & El usuario accede a la vista pacientes seleccionando la opción pacientes en el menú\\\cline{2-3} 
	 	 \multirow{2}{3.5cm}{Excepciones} & Id & Excepción \\
	 	 & 1. & Ninguna\\\cline{2-3} 
	 	 Importancia & \multicolumn{2}{p{10.25cm}}{Alta}\\
}

\tablaSmallSinColores{CU-4.1: Añadir pacientes.}{p{3cm} p{.75cm} p{9cm}}{tablaCU4.1}{
	\multicolumn{3}{p{10.25cm}}{CU-4.1: Añadir pacientes}\\}
{
	 	 Descripción &\multicolumn{2}{p{10.25cm}}{Los usuarios podrán añadir pacientes} \\\cline{2-3} 
	 	 Precondiciones & \multicolumn{2}{p{10.25cm}}{Usuario con sesión iniciada}\\\cline{2-3} 
	 	 Requisitos asociados & \multicolumn{2}{p{10.25cm}}{RF-4.1}\\\cline{2-3} 
	 	 Usuario & \multicolumn{2}{p{10.25cm}}{Administrador y oncólogo}\\\cline{2-3} 
	 	 \multirow{3}{3.5cm}{Acciones}  & Paso & Acción \\
	 	 & 1. & El usuario selecciona añadir nuevo paciente en el menú\\
	 	 & 2. & Rellena los campos con los datos del paciente\\
	 	 &3.  & Confirma el formulario\\\cline{2-3} 
	 	 \multirow{2}{3.5cm}{Excepciones} & Id & Excepción \\
	 	 & 1. & Campo en blanco\\ 
	 	 & 2. & Formato de fecha incorrecto\\\cline{2-3} 
	 	 Importancia & \multicolumn{2}{p{10.25cm}}{Alta}\\
}

\tablaSmallSinColores{CU-4.2: Añadir datos pacientes.}{p{3cm} p{.75cm} p{9cm}}{tablaCU4.2}{
	\multicolumn{3}{p{10.25cm}}{CU-4.2: Añadir datos pacientes}\\}
{
	 	 Descripción &\multicolumn{2}{p{10.25cm}}{Los usuarios podrán añadir datos a un paciente} \\\cline{2-3} 
	 	 Precondiciones & \multicolumn{2}{p{10.25cm}}{Usuario con sesión iniciada y paciente creado}\\\cline{2-3} 
	 	 Requisitos asociados & \multicolumn{2}{p{10.25cm}}{RF-4.2}\\\cline{2-3} 
	 	 Usuario & \multicolumn{2}{p{10.25cm}}{Administrador y oncólogo}\\\cline{2-3} 
	 	 \multirow{3}{3.5cm}{Acciones}  & Paso & Acción \\
	 	 & 1. & El usuario selecciona el paciente\\\cline{2-3} 
	 	 \multirow{2}{3.5cm}{Excepciones} & Id & Excepción \\
	 	 & 1. & Ninguna\\\cline{2-3} 
	 	 Importancia & \multicolumn{2}{p{10.25cm}}{Alta}\\
}

\tablaSmallSinColores{CU-4.2.1: Añadir antecedentes.}{p{3cm} p{.75cm} p{9cm}}{tablaCU4.2.1}{
	\multicolumn{3}{p{10.25cm}}{CU-4.2.1: Añadir antecedentes}\\}
{
	 	 Descripción &\multicolumn{2}{p{10.25cm}}{Los usuarios podrán añadir datos de una enfermedad a un paciente} \\\cline{2-3} 
	 	 Precondiciones & \multicolumn{2}{p{10.25cm}}{Usuario con sesión iniciada y paciente creado}\\\cline{2-3} 
	 	 Requisitos asociados & \multicolumn{2}{p{10.25cm}}{RF-4.2.1}\\\cline{2-3} 
	 	 Usuario & \multicolumn{2}{p{10.25cm}}{Administrador y oncólogo}\\\cline{2-3} 
	 	 \multirow{3}{3.5cm}{Acciones}  & Paso & Acción \\ 
	 	 & 1. & El usuario selecciona el paciente\\ 
	 	 & 2. & Se selecciona enfermedad en el menú\\ 
	 	 & 3. & Se rellena los campos con los datos de la enfermedad\\ 
	 	 & 4. & Confirma el formulario\\\cline{2-3} 
	 	 \multirow{2}{3.5cm}{Excepciones} & Id & Excepción \\ 
	 	 & 1. & Campo en blanco\\ 
	 	 & 2. & Formato de fecha incorrecto\\\cline{2-3} 
	 	 Importancia & \multicolumn{2}{p{10.25cm}}{Alta}\\
}

\tablaSmallSinColores{CU-4.2.2: Añadir antecedentes.}{p{3cm} p{.75cm} p{9cm}}{tablaCU4.2.2}{
	\multicolumn{3}{p{10.25cm}}{CU-4.2.2: Añadir antecedentes}\\}
{
	 	 Descripción &\multicolumn{2}{p{10.25cm}}{Los usuarios podrán añadir antecedentes a un paciente} \\\cline{2-3} 
	 	 Precondiciones & \multicolumn{2}{p{10.25cm}}{Usuario con sesión iniciada y paciente creado}\\\cline{2-3} 
	 	 Requisitos asociados & \multicolumn{2}{p{10.25cm}}{RF-4.2.2}\\\cline{2-3} 
	 	 Usuario & \multicolumn{2}{p{10.25cm}}{Administrador y oncólogo}\\\cline{2-3} 
	 	 \multirow{3}{3.5cm}{Acciones}  & Paso & Acción \\ 
	 	 & 1. & El usuario selecciona el paciente\\
	 	 & 2. & Se selecciona antecedentes en el menú\\
	 	 & 3. & Se rellena los campos con los datos del antecedente\\
	 	 & 4. & Confirma el formulario\\\cline{2-3} 
	 	 \multirow{2}{3.5cm}{Excepciones} & Id & Excepción \\
	 	 & 1. & Campo en blanco\\\cline{2-3} 
	 	 Importancia & \multicolumn{2}{p{10.25cm}}{Alta}\\
}

\tablaSmallSinColores{CU-4.2.3: Añadir tratamientos.}{p{3cm} p{.75cm} p{9cm}}{tablaCU4.2.3}{
	\multicolumn{3}{p{10.25cm}}{CU-4.2.3: Añadir tratamientos}\\}
{
	 	 Descripción &\multicolumn{2}{p{10.25cm}}{Los usuarios podrán añadir tratamientos a un paciente} \\\cline{2-3} 
	 	 Precondiciones & \multicolumn{2}{p{10.25cm}}{Usuario con sesión iniciada y paciente creado}\\\cline{2-3} 
	 	 Requisitos asociados & \multicolumn{2}{p{10.25cm}}{RF-4.2.3}\\\cline{2-3} 
	 	 Usuario & \multicolumn{2}{p{10.25cm}}{Administrador y oncólogo}\\\cline{2-3} 
	 	 \multirow{3}{3.5cm}{Acciones}  & Paso & Acción \\
	 	 & 1. & El usuario selecciona el paciente\\ 
	 	 & 2. & Se selecciona tratamientos en el menú\\
	 	 & 3. & Se rellena los campos con los datos del tratamiento\\
	 	 & 4. & Confirma el formulario\\\cline{2-3} 
	 	 \multirow{2}{3.5cm}{Excepciones} & Id & Excepción \\
	 	 & 1. & Campo en blanco\\
	 	 & 2. & Formato de fecha incorrecto\\\cline{2-3} 
	 	 Importancia & \multicolumn{2}{p{10.25cm}}{Alta}\\
}

\tablaSmallSinColores{CU-4.2.4: Añadir reevaluaciones.}{p{3cm} p{.75cm} p{9cm}}{tablaCU4.2.4}{
	\multicolumn{3}{p{10.25cm}}{CU-4.2.4: Añadir reevaluaciones}\\}
{
	 	 Descripción &\multicolumn{2}{p{10.25cm}}{Los usuarios podrán añadir reevaluaciones a un paciente} \\\cline{2-3} 
	 	 Precondiciones & \multicolumn{2}{p{10.25cm}}{Usuario con sesión iniciada y paciente creado}\\\cline{2-3} 
	 	 Requisitos asociados & \multicolumn{2}{p{10.25cm}}{RF-4.2.4}\\\cline{2-3} 
	 	 Usuario & \multicolumn{2}{p{10.25cm}}{Administrador y oncólogo}\\\cline{2-3} 
	 	 \multirow{3}{3.5cm}{Acciones}  & Paso & Acción \\
	 	 & 1. & El usuario selecciona el paciente\\
	 	 & 2. & Se selecciona reevaluaciones en el menú\\
	 	 & 3. & Se rellena los campos con los datos de la reevaluacion\\
	 	 & 4. & Confirma el formulario\\\cline{2-3} 
	 	 \multirow{2}{3.5cm}{Excepciones} & Id & Excepción \\
	 	 & 1. & Campo en blanco\\
	 	 & 2. & Formato de fecha incorrecto\\\cline{2-3} 
	 	 Importancia & \multicolumn{2}{p{10.25cm}}{Alta}\\
}

\tablaSmallSinColores{CU-4.2.5: Añadir seguimientos.}{p{3cm} p{.75cm} p{9cm}}{tablaCU4.2.5}{
	\multicolumn{3}{p{10.25cm}}{CU-4.2.5: Añadir seguimientos}\\}
{
	 	 Descripción &\multicolumn{2}{p{10.25cm}}{Los usuarios podrán añadir seguimientos a un paciente} \\\cline{2-3} 
	 	 Precondiciones & \multicolumn{2}{p{10.25cm}}{Usuario con sesión iniciada y paciente creado}\\\cline{2-3} 
	 	 Requisitos asociados & \multicolumn{2}{p{10.25cm}}{RF-4.2.5}\\\cline{2-3} 
	 	 Usuario & \multicolumn{2}{p{10.25cm}}{Administrador y oncólogo}\\\cline{2-3} 
	 	 \multirow{3}{3.5cm}{Acciones}  & Paso & Acción \\
	 	 & 1. & El usuario selecciona el paciente\\
	 	 & 2. & Se selecciona seguimientos en el menú\\
	 	 & 3. & Se rellena los campos con los datos del seguimiento\\
	 	 & 4. & Confirma el formulario\\\cline{2-3} 
	 	 \multirow{2}{3.5cm}{Excepciones} & Id & Excepción \\
	 	 & 1. & Campo en blanco\\
	 	 & 2. & Formato de fecha incorrecto\\\cline{2-3} 
	 	 Importancia & \multicolumn{2}{p{10.25cm}}{Alta}\\
}

\tablaSmallSinColores{CU-4.2.6: Añadir comentarios.}{p{3cm} p{.75cm} p{9cm}}{tablaCU4.2.6}{
	\multicolumn{3}{p{10.25cm}}{CU-4.2.6: Añadir comentarios}\\}
{
	 	 Descripción &\multicolumn{2}{p{10.25cm}}{Los usuarios podrán añadir comentarios a un paciente} \\\cline{2-3} 
	 	 Precondiciones & \multicolumn{2}{p{10.25cm}}{Usuario con sesión iniciada y paciente creado}\\\cline{2-3} 
	 	 Requisitos asociados & \multicolumn{2}{p{10.25cm}}{RF-4.2.6}\\\cline{2-3} 
	 	 Usuario & \multicolumn{2}{p{10.25cm}}{Administrador y oncólogo}\\\cline{2-3} 
	 	 \multirow{3}{3.5cm}{Acciones}  & Paso & Acción \\
	 	 & 1. & El usuario selecciona el paciente\\
	 	 & 2. & Se selecciona comentarios en el menú\\
	 	 & 3. & Se rellena el campo con el comentario\\
	 	 & 4. & Confirma el formulario\\\cline{2-3} 
	 	 \multirow{2}{3.5cm}{Excepciones} & Id & Excepción \\
	 	 & 1. & Campo en blanco\\\cline{2-3} 
	 	 Importancia & \multicolumn{2}{p{10.25cm}}{Media}\\
}

\tablaSmallSinColores{CU-4.3: Modificar pacientes.}{p{3cm} p{.75cm} p{9cm}}{tablaCU4.3}{
	\multicolumn{3}{p{10.25cm}}{CU-4.3: Modificar pacientes}\\}
{
	 	 Descripción &\multicolumn{2}{p{10.25cm}}{Los usuarios podrán modificar los datos de los pacientes} \\\cline{2-3} 
	 	 Precondiciones & \multicolumn{2}{p{10.25cm}}{Usuario con sesión iniciada y paciente creado}\\\cline{2-3} 
	 	 Requisitos asociados & \multicolumn{2}{p{10.25cm}}{RF-4.3}\\\cline{2-3} 
	 	 Usuario & \multicolumn{2}{p{10.25cm}}{Administrador y oncólogo}\\\cline{2-3} 
	 	 \multirow{3}{3.5cm}{Acciones}  & Paso & Acción \\
	 	 & 1. & El usuario selecciona el paciente\\\cline{2-3} 
	 	 \multirow{2}{3.5cm}{Excepciones} & Id & Excepción \\
	 	 & 1 & Ninguna\\\cline{2-3} 
	 	 Importancia & \multicolumn{2}{p{10.25cm}}{Alta}\\
}

\tablaSmallSinColores{CU-4.3.1: Modificar datos demográficos.}{p{3cm} p{.75cm} p{9cm}}{tablaCU4.3.1}{
	\multicolumn{3}{p{10.25cm}}{CU-4.3.1: Modificar datos demográficos}\\}
{
	 	 Descripción &\multicolumn{2}{p{10.25cm}}{Los usuarios podrán modificar los datos de demográficos de un paciente} \\\cline{2-3} 
	 	 Precondiciones & \multicolumn{2}{p{10.25cm}}{Usuario con sesión iniciada y paciente creado}\\\cline{2-3} 
	 	 Requisitos asociados & \multicolumn{2}{p{10.25cm}}{RF-4.3.1}\\\cline{2-3} 
	 	 Usuario & \multicolumn{2}{p{10.25cm}}{Administrador y oncólogo}\\\cline{2-3} 
	 	 \multirow{3}{3.5cm}{Acciones}  & Paso & Acción \\
	 	 & 1. & El usuario selecciona el paciente\\
	 	 & 2. & Se modifican los campos que se van a cambiar\\
	 	 & 3. & Confirma el formulario\\\cline{2-3} 
	 	 \multirow{2}{3.5cm}{Excepciones} & Id & Excepción \\
	 	 & 1. & Campo en blanco\\
	 	 & 2. & Formato de fecha incorrecto\\\cline{2-3} 
	 	 Importancia & \multicolumn{2}{p{10.25cm}}{Alta}\\
}

\tablaSmallSinColores{CU-4.3.2: Modificar datos enfermedad.}{p{3cm} p{.75cm} p{9cm}}{tablaCU4.3.2}{
	\multicolumn{3}{p{10.25cm}}{CU-4.3.2: Modificar datos enfermedad}\\}
{
	 	 Descripción &\multicolumn{2}{p{10.25cm}}{Los usuarios podrán modificar los datos de la enfermedad de un paciente} \\\cline{2-3} 
	 	 Precondiciones & \multicolumn{2}{p{10.25cm}}{Usuario con sesión iniciada y paciente creado}\\\cline{2-3} 
	 	 Requisitos asociados & \multicolumn{2}{p{10.25cm}}{RF-4.3.2}\\\cline{2-3} 
	 	 Usuario & \multicolumn{2}{p{10.25cm}}{Administrador y oncólogo}\\\cline{2-3} 
	 	 \multirow{3}{3.5cm}{Acciones}  & Paso & Acción \\
	 	 & 1. & El usuario selecciona el paciente\\
	 	 & 2. & El usuario selecciona enfermedad en el menú\\
	 	 & 3. & Se modifican los campos que se van a cambiar\\
	 	 & 4. & Confirma el formulario\\\cline{2-3} 
	 	 \multirow{2}{3.5cm}{Excepciones} & Id & Excepción \\
	 	 & 1. & Campo en blanco\\
	 	 & 2. & Formato de fecha incorrecto\\\cline{2-3} 
	 	 Importancia & \multicolumn{2}{p{10.25cm}}{Alta}\\
}

\tablaSmallSinColores{CU-4.3.3: Modificar datos antecedentes.}{p{3cm} p{.75cm} p{9cm}}{tablaCU4.3.3}{
	\multicolumn{3}{p{10.25cm}}{CU-4.3.3: Modificar datos antecedentes}\\}
{
	 	 Descripción &\multicolumn{2}{p{10.25cm}}{Los usuarios podrán modificar los datos de los antecedentes de un paciente} \\\cline{2-3} 
	 	 Precondiciones & \multicolumn{2}{p{10.25cm}}{Usuario con sesión iniciada y paciente creado}\\\cline{2-3} 
	 	 Requisitos asociados & \multicolumn{2}{p{10.25cm}}{RF-4.3.3}\\\cline{2-3} 
	 	 Usuario & \multicolumn{2}{p{10.25cm}}{Administrador y oncólogo}\\\cline{2-3} 
	 	 \multirow{3}{3.5cm}{Acciones}  & Paso & Acción \\
	 	 & 1. & El usuario selecciona el paciente\\
	 	 & 2. & El usuario selecciona antecedentes en el menú\\
	 	 & 3. & Se modifican los campos que se van a cambiar\\
	 	 & 4. & Confirma el formulario\\\cline{2-3} 
	 	 \multirow{2}{3.5cm}{Excepciones} & Id & Excepción \\
	 	 & 1. & Campo en blanco\\\cline{2-3} 
	 	 Importancia & \multicolumn{2}{p{10.25cm}}{Alta}\\
}

\tablaSmallSinColores{CU-4.3.4: Modificar datos tratamientos.}{p{3cm} p{.75cm} p{9cm}}{tablaCU4.3.4}{
	\multicolumn{3}{p{10.25cm}}{CU-4.3.4: Modificar datos tratamientos}\\}
{
	 	 Descripción &\multicolumn{2}{p{10.25cm}}{Los usuarios podrán modificar los datos de los tratamientos de un paciente} \\\cline{2-3} 
	 	 Precondiciones & \multicolumn{2}{p{10.25cm}}{Usuario con sesión iniciada y paciente creado}\\\cline{2-3} 
	 	 Requisitos asociados & \multicolumn{2}{p{10.25cm}}{RF-4.3.4}\\\cline{2-3} 
	 	 Usuario & \multicolumn{2}{p{10.25cm}}{Administrador y oncólogo}\\\cline{2-3} 
	 	 \multirow{3}{3.5cm}{Acciones}  & Paso & Acción \\
	 	 & 1. & El usuario selecciona el usuario\\
	 	 & 2. & El usuario selecciona tratamientos en el menú\\
	 	 & 3. & Se modifican los campos que se van a cambiar\\
	 	 & 4. & Confirma el formulario\\\cline{2-3} 
	 	 \multirow{2}{3.5cm}{Excepciones} & Id & Excepción \\
	 	 & 1. & Campo en blanco\\
	 	 & 2. & Formato de fecha incorrecto\\\cline{2-3} 
	 	 Importancia & \multicolumn{2}{p{10.25cm}}{Alta}\\
}

\tablaSmallSinColores{CU-4.3.5: Modificar datos reevaluaciones.}{p{3cm} p{.75cm} p{9cm}}{tablaCU4.3.5}{
	\multicolumn{3}{p{10.25cm}}{CU-4.3.5: Modificar datos reevaluaciones}\\}
{
	 	 Descripción &\multicolumn{2}{p{10.25cm}}{Los usuarios podrán modificar los datos de las reevaluaciones de un paciente} \\\cline{2-3} 
	 	 Precondiciones & \multicolumn{2}{p{10.25cm}}{Usuario con sesión iniciada y paciente creado}\\\cline{2-3} 
	 	 Requisitos asociados & \multicolumn{2}{p{10.25cm}}{RF-4.3.5}\\\cline{2-3} 
	 	 Usuario & \multicolumn{2}{p{10.25cm}}{Administrador y oncólogo}\\\cline{2-3} 
	 	 \multirow{3}{3.5cm}{Acciones}  & Paso & Acción \\
	 	 & 1. & El usuario selecciona el usuario\\
	 	 & 2. & El usuario selecciona reevaluaciones en el menú\\
	 	 & 3. & Se modifican los campos que se van a cambiar\\
	 	 & 4. & Confirma el formulario\\\cline{2-3} 
	 	 \multirow{2}{3.5cm}{Excepciones} & Id & Excepción \\
	 	 & 1. & Campo en blanco\\
	 	 & 2. & Formato de fecha incorrecto\\\cline{2-3} 
	 	 Importancia & \multicolumn{2}{p{10.25cm}}{Alta}\\
}

\tablaSmallSinColores{CU-4.3.6: Modificar datos seguimientos.}{p{3cm} p{.75cm} p{9cm}}{tablaCU4.3.6}{
	\multicolumn{3}{p{10.25cm}}{CU-4.3.6: Modificar datos seguimientos}\\}
{
	 	 Descripción &\multicolumn{2}{p{10.25cm}}{Los usuarios podrán modificar los datos de los seguimientos de un paciente} \\\cline{2-3} 
	 	 Precondiciones & \multicolumn{2}{p{10.25cm}}{Usuario con sesión iniciada y paciente creado}\\\cline{2-3} 
	 	 Requisitos asociados & \multicolumn{2}{p{10.25cm}}{RF-4.3.6}\\\cline{2-3} 
	 	 Usuario & \multicolumn{2}{p{10.25cm}}{Administrador y oncólogo}\\\cline{2-3} 
	 	 \multirow{3}{3.5cm}{Acciones}  & Paso & Acción \\
	 	 & 1. & El usuario selecciona el usuario\\
	 	 & 2. & El usuario selecciona seguimientos en el menú\\
	 	 & 3. & Se modifican los campos que se van a cambiar\\
	 	 & 4. & Confirma el formulario\\\cline{2-3} 
	 	 \multirow{2}{3.5cm}{Excepciones} & Id & Excepción \\
	 	 & 1. & Campo en blanco\\
	 	 & 2. & Formato de fecha incorrecto\\\cline{2-3} 
	 	 Importancia & \multicolumn{2}{p{10.25cm}}{Alta}\\
}

\tablaSmallSinColores{CU-4.3.7: Modificar comentarios.}{p{3cm} p{.75cm} p{9cm}}{tablaCU4.3.7}{
	\multicolumn{3}{p{10.25cm}}{CU-4.3.7: Modificar comentarios}\\}
{
	 	 Descripción &\multicolumn{2}{p{10.25cm}}{Los usuarios podrán modificar los comentarios a un paciente} \\\cline{2-3} 
	 	 Precondiciones & \multicolumn{2}{p{10.25cm}}{Usuario con sesión iniciada y paciente creado}\\\cline{2-3} 
	 	 Requisitos asociados & \multicolumn{2}{p{10.25cm}}{RF-4.3.7}\\\cline{2-3} 
	 	 Usuario & \multicolumn{2}{p{10.25cm}}{Administrador y oncólogo}\\\cline{2-3} 
	 	 \multirow{3}{3.5cm}{Acciones}  & Paso & Acción \\
	 	 & 1. & El usuario selecciona el usuario\\
	 	 & 2. & El usuario selecciona comentarios en el menú\\
	 	 & 3. & Se modifican el comentario\\
	 	 & 4. & Confirma el formulario\\\cline{2-3} 
	 	 \multirow{2}{3.5cm}{Excepciones} & Id & Excepción \\
	 	 & 1. & Campo en blanco\\\cline{2-3} 
	 	 Importancia & \multicolumn{2}{p{10.25cm}}{Alta}\\
}

\tablaSmallSinColores{CU-4.4: Eliminar paciente.}{p{3cm} p{.75cm} p{9cm}}{tablaCU4.4}{
	\multicolumn{3}{p{10.25cm}}{CU-4.4: Eliminar paciente}\\}
{
	 	 Descripción &\multicolumn{2}{p{10.25cm}}{Los usuarios podrán eliminar pacientes} \\\cline{2-3} 
	 	 Precondiciones & \multicolumn{2}{p{10.25cm}}{Usuario con sesión iniciada y paciente creado}\\\cline{2-3} 
	 	 Requisitos asociados & \multicolumn{2}{p{10.25cm}}{RF-4.4}\\\cline{2-3} 
	 	 Usuario & \multicolumn{2}{p{10.25cm}}{Administrador y oncólogo}\\\cline{2-3} 
	 	 \multirow{3}{3.5cm}{Acciones}  & Paso & Acción \\
	 	 & 1. & El usuario selecciona eliminar usuario\\
	 	 & 2. & El usuario selecciona el usuario a eliminar\\
	 	 & 3. & Confirma el formulario\\\cline{2-3} 
	 	 \multirow{2}{3.5cm}{Excepciones} & Id & Excepción \\
	 	 & 1. & Ninguna\\\cline{2-3} 
	 	 Importancia & \multicolumn{2}{p{10.25cm}}{Alta}\\
}

\tablaSmallSinColores{CU-5: Visualización de los datos generales.}{p{3cm} p{.75cm} p{9cm}}{tablaCU5}{
	\multicolumn{3}{p{10.25cm}}{CU-5: Visualización de los datos generales}\\}
{
	 	 Descripción &\multicolumn{2}{p{10.25cm}}{Los usuarios podrán visualizar características de los pacientes mediante gráficas} \\\cline{2-3} 
	 	 Precondiciones & \multicolumn{2}{p{10.25cm}}{Usuario con sesión iniciada y pacientes en la base de datos}\\\cline{2-3} 
	 	 Requisitos asociados & \multicolumn{2}{p{10.25cm}}{RF-5, RF-5.1}\\\cline{2-3} 
	 	 Usuario & \multicolumn{2}{p{10.25cm}}{Administrador y oncólogo}\\\cline{2-3} 
	 	 \multirow{3}{3.5cm}{Acciones}  & Paso & Acción \\
	 	 & 1. & El usuario selecciona gráficas en el menú desplegable\\
	 	 & 2. & El usuario selecciona gráficas generales en el menú lateral\\
	 	 & 3. & El usuario selecciona los campos en los que desea separar las gráficas\\
	 	 & 4. & Confirma el formulario\\\cline{2-3} 
	 	 \multirow{2}{3.5cm}{Excepciones} & Id & Excepción \\
	 	 & 1. & Que no existan pacientes con esos atributos en específico\\\cline{2-3} 
	 	 Importancia & \multicolumn{2}{p{10.25cm}}{Media}\\
}

\tablaSmallSinColores{CU-6: Visualización de los datos numéricos.}{p{3cm} p{.75cm} p{9cm}}{tablaCU6}{
	\multicolumn{3}{p{10.25cm}}{CU-6: Visualización de los datos numéricos}\\}
{
	 	 Descripción &\multicolumn{2}{p{10.25cm}}{Los usuarios podrán visualizar diferentes atributos estadísticos de los datos numéricos} \\\cline{2-3} 
	 	 Precondiciones & \multicolumn{2}{p{10.25cm}}{Usuario con sesión iniciada y pacientes en la base de datos}\\\cline{2-3} 
	 	 Requisitos asociados & \multicolumn{2}{p{10.25cm}}{RF-5.2}\\\cline{2-3} 
	 	 Usuario & \multicolumn{2}{p{10.25cm}}{Administrador y oncólogo}\\\cline{2-3} 
	 	 \multirow{3}{3.5cm}{Acciones}  & Paso & Acción \\
	 	 & 1. & El usuario selecciona gráficas en el menú desplegable\\
	 	 & 2. & El usuario selecciona percentiles en el menú lateral\\
	 	 & 3. & El usuario selecciona el campo que desea obtener información\\
	 	 & 4. & Confirma el formulario\\\cline{2-3} 
	 	 \multirow{2}{3.5cm}{Excepciones} & Id & Excepción \\
	 	 & 1. & Que no existan pacientes con ese atributos en específico\\\cline{2-3} 
	 	 Importancia & \multicolumn{2}{p{10.25cm}}{Media}\\
}

\tablaSmallSinColores{CU-7: Visualización de los datos individualmente.}{p{3cm} p{.75cm} p{9cm}}{tablaCU7}{
	\multicolumn{3}{p{10.25cm}}{CU-7: Visualización de los datos individualmente}\\}
{
	 	 Descripción &\multicolumn{2}{p{10.25cm}}{Los usuarios podrán visualizar los datos de los pacientes individualmente} \\\cline{2-3} 
	 	 Precondiciones & \multicolumn{2}{p{10.25cm}}{Usuario con sesión iniciada y pacientes en la base de datos}\\\cline{2-3} 
	 	 Requisitos asociados & \multicolumn{2}{p{10.25cm}}{RF-5.3}\\\cline{2-3} 
	 	 Usuario & \multicolumn{2}{p{10.25cm}}{Administrador y oncólogo}\\\cline{2-3} 
	 	 \multirow{3}{3.5cm}{Acciones}  & Paso & Acción \\
	 	 & 1. & El usuario selecciona pacientes en el menú\\
	 	 & 2. & El usuario selecciona el paciente del cual desea ver los datos\\\cline{2-3} 
	 	 \multirow{2}{3.5cm}{Excepciones} & Id & Excepción \\
	 	 & 1. & Ninguna\\\cline{2-3} 
	 	 Importancia & \multicolumn{2}{p{10.25cm}}{Baja}\\
}

\tablaSmallSinColores{CU-8: Exportación de los datos.}{p{3cm} p{.75cm} p{9cm}}{tablaCU8}{
	\multicolumn{3}{p{10.25cm}}{CU-8: Exportación de los datos}\\}
{
	 	 Descripción &\multicolumn{2}{p{10.25cm}}{Los usuarios podrán exportar los datos de los paciente a un fichero Excel} \\\cline{2-3} 
	 	 Precondiciones & \multicolumn{2}{p{10.25cm}}{Usuario con sesión iniciada y pacientes en la base de datos}\\\cline{2-3} 
	 	 Requisitos asociados & \multicolumn{2}{p{10.25cm}}{RF-5.4}\\\cline{2-3} 
	 	 Usuario & \multicolumn{2}{p{10.25cm}}{Administrador y oncólogo}\\\cline{2-3} 
	 	 \multirow{3}{3.5cm}{Acciones}  & Paso & Acción \\
	 	 & 1. & El usuario selecciona exportar datos en el menú\\
	 	 & 2. & El usuario rellena el campo con el nombre del fichero\\
	 	 & 3. & El usuario confirma el formulario\\\cline{2-3} 
	 	 \multirow{2}{3.5cm}{Excepciones} & Id & Excepción \\
	 	 & 1. & Campo en blanco\\\cline{2-3} 
	 	 Importancia & \multicolumn{2}{p{10.25cm}}{Alta}\\
}

\tablaSmallSinColores{CU-9: Visualización de los datos en la gráfica kaplan meier.}{p{3cm} p{.75cm} p{9cm}}{tablaCU9}{
	\multicolumn{3}{p{10.25cm}}{CU-8: Visualización de los datos en la gráfica kaplan meier}\\}
{
	 	 Descripción &\multicolumn{2}{p{10.25cm}}{Los usuarios podrán visualizar los datos de supervivencia en una grafica kaplan meier} \\\cline{2-3} 
	 	 Precondiciones & \multicolumn{2}{p{10.25cm}}{Usuario con sesión iniciada y pacientes en la base de datos}\\\cline{2-3} 
	 	 Requisitos asociados & \multicolumn{2}{p{10.25cm}}{RF-5.5}\\\cline{2-3} 
	 	 Usuario & \multicolumn{2}{p{10.25cm}}{Administrador y oncólogo}\\\cline{2-3} 
	 	 \multirow{3}{3.5cm}{Acciones}  & Paso & Acción \\
	 	 & 1. & El usuario selecciona gráficas en el menú desplegable\\
	 	 & 2. & El usuario selecciona Kaplan meier en el menú lateral\\
	 	 & 3. & El usuario rellena selecciona el campo por el que dividir la gráfica\\
	 	 & 4. & El usuario confirma el formulario\\\cline{2-3} 
	 	 \multirow{2}{3.5cm}{Excepciones} & Id & Excepción \\
	 	 & 1. & Insuficientes datos para realizar la gráfica\\\cline{2-3} 
	 	 Importancia & \multicolumn{2}{p{10.25cm}}{Alta}\\
}

\tablaSmallSinColores{CU-10: Creación de base de datos sintética.}{p{3cm} p{.75cm} p{9cm}}{tablaCU10}{
	\multicolumn{3}{p{10.25cm}}{CU-10: Creación de base de datos sintética}\\}
{
	 	 Descripción &\multicolumn{2}{p{10.25cm}}{Los usuarios podrán crear una base de datos con datos sintéticos} \\\cline{2-3} 
	 	 Precondiciones & \multicolumn{2}{p{10.25cm}}{Usuario con sesión iniciadas}\\\cline{2-3} 
	 	 Requisitos asociados & \multicolumn{2}{p{10.25cm}}{RF-6}\\\cline{2-3} 
	 	 Usuario & \multicolumn{2}{p{10.25cm}}{Administrador y oncólogo}\\\cline{2-3} 
	 	 \multirow{3}{3.5cm}{Acciones}  & Paso & Acción \\
	 	 & 1. & El usuario selecciona base sintética en el menú\\
	 	 & 2. & El usuario rellena los campos\\
	 	 & 3. & El usuario confirma el formulario\\\cline{2-3} 
	 	 \multirow{2}{3.5cm}{Excepciones} & Id & Excepción \\
	 	 & 1. & Datos introducidos fuera de los rangos permitidos\\\cline{2-3} 
	 	 Importancia & \multicolumn{2}{p{10.25cm}}{Media}\\
}

\tablaSmallSinColores{CU-11: Filtrado de pacientes tabla.}{p{3cm} p{.75cm} p{9cm}}{tablaCU11}{
	\multicolumn{3}{p{10.25cm}}{CU-11: Filtrado de pacientes tabla}\\}
{
	 	 Descripción &\multicolumn{2}{p{10.25cm}}{Los usuarios podrán filtrar los pacientes por atributos y ver la tabla de gestión de pacientes con los pacientes que cumplan este filtro} \\\cline{2-3} 
	 	 Precondiciones & \multicolumn{2}{p{10.25cm}}{Usuario con sesión iniciadas}\\\cline{2-3} 
	 	 Requisitos asociados & \multicolumn{2}{p{10.25cm}}{RF-7.1}\\\cline{2-3} 
	 	 Usuario & \multicolumn{2}{p{10.25cm}}{Administrador y oncólogo}\\\cline{2-3} 
	 	 \multirow{3}{3.5cm}{Acciones}  & Paso & Acción \\
	 	 & 1. & El usuario selecciona paciente en el menú\\
	 	 & 2. & El usuario da clic en el botón Aplicar filtro\\
	 	 & 3. & El usuario selecciona los atributos por los que realizar el filtro\\
	 	 & 4. & El usuario da clic en el botón realizar filtrado\\\cline{2-3} 
	 	 \multirow{2}{3.5cm}{Excepciones} & Id & Excepción \\
	 	 & 1. & Datos introducidos fuera de los rangos permitidos\\\cline{2-3} 
	 	 Importancia & \multicolumn{2}{p{10.25cm}}{Media}\\
}

\tablaSmallSinColores{CU-12: Filtrado de pacientes y visualización de datos.}{p{3cm} p{.75cm} p{9cm}}{tablaCU12}{
	\multicolumn{3}{p{10.25cm}}{CU-12: Filtrado de pacientes y visualización de datos}\\}
{
	 	 Descripción &\multicolumn{2}{p{10.25cm}}{Los usuarios podrán filtrar los pacientes por atributos y visualizar características de los pacientes mediante gráficas} \\\cline{2-3} 
	 	 Precondiciones & \multicolumn{2}{p{10.25cm}}{Usuario con sesión iniciadas}\\\cline{2-3} 
	 	 Requisitos asociados & \multicolumn{2}{p{10.25cm}}{RF-7.2}\\\cline{2-3} 
	 	 Usuario & \multicolumn{2}{p{10.25cm}}{Administrador y oncólogo}\\\cline{2-3} 
	 	 \multirow{3}{3.5cm}{Acciones}  & Paso & Acción \\
	 	 & 1. & El usuario selecciona gráficas en el menú\\
	 	 & 2. & El usuario selecciona gráficas generales en el menú lateral\\
	 	 & 3. & El usuario da clic en el botón Aplicar filtro\\
	 	 & 4. & El usuario selecciona los atributos por los que realizar el filtro\\
	 	 & 5. & El usuario da clic en el botón realizar filtrado\\
	 	 & 6. & El usuario selecciona los campos en los que desea separar las gráficas\\
	 	 & 7. & El usuario confirma el formulario\\\cline{2-3} 
	 	 \multirow{2}{3.5cm}{Excepciones} & Id & Excepción \\
	 	 & 1. & No existe ningún paciente con ese filtro\\
	 	 & 2. & No se selecciona ningún filtro\\\cline{2-3} 
	 	 Importancia & \multicolumn{2}{p{10.25cm}}{Media}\\
}

\tablaSmallSinColores{CU-13: Filtrado de pacientes y visualización de datos numéricos.}{p{3cm} p{.75cm} p{9cm}}{tablaCU13}{
	\multicolumn{3}{p{10.25cm}}{CU-13: Filtrado de pacientes y visualización de datos numéricos}\\}
{
	 	 Descripción &\multicolumn{2}{p{10.25cm}}{Los usuarios podrán filtrar los pacientes por atributos y visualizar datos estadísticos de los datos numéricos} \\\cline{2-3} 
	 	 Precondiciones & \multicolumn{2}{p{10.25cm}}{Usuario con sesión iniciadas}\\\cline{2-3} 
	 	 Requisitos asociados & \multicolumn{2}{p{10.25cm}}{RF-7.3}\\\cline{2-3} 
	 	 Usuario & \multicolumn{2}{p{10.25cm}}{Administrador y oncólogo}\\\cline{2-3} 
	 	 \multirow{3}{3.5cm}{Acciones}  & Paso & Acción \\
	 	 & 1. & El usuario selecciona gráficas en el menú\\
	 	 & 2. & El usuario selecciona percentiles en el menú lateral\\
	 	 & 3. & El usuario da clic en el botón Aplicar filtro\\
	 	 & 4. & El usuario selecciona los atributos por los que realizar el filtro\\
	 	 & 5. & El usuario da clic en el botón realizar filtrado\\
	 	 & 6. & El usuario selecciona el campo que desea obtener información\\
	 	 & 7. & El usuario confirma el formulario\\\cline{2-3} 
	 	 \multirow{2}{3.5cm}{Excepciones} & Id & Excepción \\
	 	 & 1. & No existe ningún paciente con ese filtro\\
	 	 & 2. & No se selecciona ningún filtro\\\cline{2-3} 
	 	 Importancia & \multicolumn{2}{p{10.25cm}}{Media}\\
}

\tablaSmallSinColores{CU-14: Filtrado de pacientes y visualización de la gráfica kaplan meier.}{p{3cm} p{.75cm} p{9cm}}{tablaCU14}{
	\multicolumn{3}{p{10.25cm}}{CU-14: Filtrado de pacientes y visualización de la gráfica kaplan meier}\\}
{
	 	 Descripción &\multicolumn{2}{p{10.25cm}}{Los usuarios podrán filtrar los pacientes por atributos y visualizar la tabla de supervivencia kaplan meier} \\\cline{2-3} 
	 	 Precondiciones & \multicolumn{2}{p{10.25cm}}{Usuario con sesión iniciadas}\\\cline{2-3} 
	 	 Requisitos asociados & \multicolumn{2}{p{10.25cm}}{RF-7.3}\\\cline{2-3} 
	 	 Usuario & \multicolumn{2}{p{10.25cm}}{Administrador y oncólogo}\\\cline{2-3} 
	 	 \multirow{3}{3.5cm}{Acciones}  & Paso & Acción \\
	 	 & 1. & El usuario selecciona gráficas en el menú\\
	 	 & 2. & El usuario selecciona kaplan meier en el menú lateral\\
	 	 & 3. & El usuario da clic en el botón Aplicar filtro\\
	 	 & 4. & El usuario selecciona los atributos por los que realizar el filtro\\
	 	 & 5. & El usuario da clic en el botón realizar filtrado\\
	 	 & 6. & El usuario selecciona el campo por el que desea dividir la gráfica\\
	 	 & 7. & El usuario confirma el formulario\\\cline{2-3} 
	 	 \multirow{2}{3.5cm}{Excepciones} & Id & Excepción \\
	 	 & 1. & No existe ningún paciente con ese filtro\\
	 	 & 2. & No se selecciona ningún filtro\\\cline{2-3} 
	 	 Importancia & \multicolumn{2}{p{10.25cm}}{Media}\\
}
\apendice{Especificación de diseño}

\section{Introducción}

\section{Diseño de datos}

\section{Diseño procedimental}

\section{Diseño arquitectónico}



\apendice{Documentación técnica de programación}

\section{Introducción}

En este anexo se va explicar la estructura del proyecto, los requisitos necesarios para poder ejecutarlo desde cualquier ordenador  

\section{Estructura de directorios}

El contenido del repositorio se estructura principalmente en 4 directorios:

\begin{itemize}
    \item \textbf{/Base de datos sintética:} Contiene el script para generar la base de datos sintética.
    \item \textbf{/Lavarel:} Contiene el proyecto de la aplicación web junto con los test, las migraciones, los seeders y el enrutado.
    \item \textbf{/SQL:} Contiene los scripts SQL que crean la estructura de la base de datos e introducen los datos necesarios para que funcione la aplicación.
    \item \textbf{/docs:} Contiene la documentación del proyecto. 
\end{itemize}

Dentro del proyecto Laravel los directorios más importantes son:

\begin{itemize}
    \item \textbf{/app:} Contiene el script para generar la base de datos sintética.
    \begin{itemize}
        \item \textbf{/Exports:} Contiene los ficheros necesarios para realizar la exportación de la base de datos a el fichero Excel.
        \item \textbf{/Http:} 
        \begin{itemize}
            \item \textbf{/Controllers:}
            \item \textbf{/Middleware:}
        \end{itemize}
        \item \textbf{/JoinsTablas:} Contiene el fichero que permite devuelve un outter join entre todas las tablas.
        \item \textbf{/Models:} Contiene los modelos de las tablas de la base de datos, cada uno con todos sus atributos y con sus relaciones.
        \item \textbf{/Rules:} Contiene las reglas que hemos añadido para la validación de datos como por ejemplo para comprobar correo único en la base de datos.
        \item \textbf{/Utilidades:} Contiene el fichero que permite tanto la encriptación de datos como la desencriptación .
    \end{itemize}
    \item \textbf{/config:} Contiene los ficheros de configuración del proyecto Laravel.
    \begin{itemize}
        \item \textbf{/app.php:} Fichero de configuración generales del proyecto, como por ejemplo nombre de la aplicación, zona horaria, etc.
        \item \textbf{/auth.php} Fichero de configuración de la autentificación del login.
    \end{itemize}
    \item \textbf{/database:} Contiene los ficheros para poner la base de datos a punto para poder usar la aplicación.
    \begin{itemize}
        \item \textbf{/migrations:} Contiene ficheros para poder crear la estructura de cada una de las tablas, con sus atributos y sus relaciones.
        \item \textbf{/seeders:} Contiene ficheros para introducir los datos necesarios para que funcione la aplicación.
    \end{itemize}
    \item \textbf{/public:} 
    \begin{itemize}
        \item \textbf{/css:} Contiene los ficheros \textit{.css} que van a servir para dar formato a nuestra web.
        \item \textbf{/js:} Contiene los ficheros \textit{.js} los cuales va a añadir animaciones y efectos a nuestra web.
    \end{itemize}
    \item \textbf{/resources/views:} Contiene los ficheros HTML de todas las vistas de la web.
    \item \textbf{/routes/web.php:} Contiene todas las rutas de la web con sus correspondientes peticiones HTTP-
    \item \textbf{/tests:} Contiene todos los test realizados con phpunit.
    \begin{itemize}
        \item \textbf{/Featute:} Contiene los test realizados de la gestión de pacientes y usuarios.
        \item \textbf{/Ini:} Contiene ficheros para crear pacientes y usuarios para los tests.
    \end{itemize}
    \item \textbf{/.env:} Contiene las variables de entorno.
    \item \textbf{/phpunit.xml:} Contiene la configuración para poder ejecutar los tests.
\end{itemize}

\section{Manual del programador} \label{manual}

En este apartado se van a comentar todos los programas y versiones que van a ser necesarias para correr la aplicación en local.

\subsection{WAMP}

Es un conjunto de programas para permitir el desarrollo de aplicaciones web en local en sistemas Windows, esta consta de 4 herramientas principales:
\begin{itemize}
    \item \textbf{Apache:} Servidor web HTTP de código abierto.
    \item \textbf{MySQL:} Base de datos relacional 
    \item \textbf{PHP:} Lenguaje de programación centrado en el desarrollo web.
    \item \textbf{PhpMyAdmin:} Programa para gestionar bases de datos.
\end{itemize}

WAMP se descarga en la web \href{https://www.wampserver.com/}{https://www.wampserver.com/}, en el instalador usar todos los valores por defecto y la versión que se esta usando es la 3.2.3.3.

\subsection{composer}

Es un gestor de paquetes para PHP y es el que en nuestro proyecto nos ha permitido tanto descargar Laravel como las librerías de PHP que hemos necesitado. La versión que se esta usando es la 2.0.13.

Composer se descarga en la web \href{https://getcomposer.org/}{https://getcomposer.org/}, en el instalador usar todos los valores por defecto, excepto cuando pregunta por la version de PHP que queremos usar, en este caso seleccionaremos "\textit{C:\textbackslash{}wamp64\textbackslash{}bin\textbackslash{}php\\\textbackslash{}php7.4.9\textbackslash{}php.exe}" como podemos ver en la imagen \ref{fig:composer}.

\imagenTam{composer}{Versión de PHP a elegir en el instalador de composer}{0.7}

\subsection{Python 3}

Es un lenguaje de programación, el cual se ha usado para la creación de la base de datos sintética, y es necesario tener Python instalado en el ordenador para permitir correr este script. La versión que se esta usando es la 3.9.5.

Python se descarga en la web \href{https://www.python.org/downloads/release/python-395/}{https://www.python.org/downloads/release\\/python-395/}, en el instalador seleccionar las opciones por defecto, excepto en la localización que vamos a poner "\textit{C:\textbackslash{}Python3.9\textbackslash{}}" como podemos ver en la imagen \ref{fig:python}.

\newpage

\subsubsection{Librerías Python}

Para este proyecto va a ser necesario tener instaladas la siguiente librerías de Python:
\begin{itemize}
    \item numpy
    \item mysql.connector
    \item faker
    \item datetime
\end{itemize}

Para instalar las siguiente librerías podemos usar el comando "\textit{pip install nombre de la librería}"

\imagenTam{python}{Localización a elegir en el instalador de python}{0.7}

\section{Compilación, instalación y ejecución del proyecto}

En este apartado se va a explicar como instalar y ejecutar el proyecto en un entorno local.

\subsection{Importar el proyecto}

Una vez tengamos instalado todos los programas y librerías como se ha explicado en el punto \ref{manual}, lo primero vamos a hacer va a ser importar el proyecto, para esto seguiremos los siguientes pasos:

\begin{enumerate}
    \item Acceder al repositorio: \href{https://github.com/CarlosOrtu/TFG_HUBU}{https://github.com/CarlosOrtu/TFG\_HUBU}.
    \item Descargar el contenido del repositorio desde \textbf{Code > Download ZIP} como se puede ver en la imagen \ref{fig:git_descargar}
    \imagenTam{git_descargar}{Descargar el proyecto GitHub}{0.5}
    \item Descomprimir el proyecto.
\end{enumerate}
    
\subsection{Añadir el proyecto a WAMP}

Una vez tenemos todos los archivos del proyecto, lo siguiente que tendremos que hacer será añadir este proyecto a WAMP. Para esto seguiremos los siguientes pasos.

\begin{enumerate}
    \item Accedemos a la ruta "\textit{C:\textbackslash{}wamp64\textbackslash{}www}"{} o a la ruta donde se haya descargado WAMP y posteriormente a la carpeta www.
    \item Creamos una carpeta con el nombre TFG\_HUBU.
    \item Pegamos todos los ficheros que hemos descomprimido del proyecto dentro de esta carpeta.
\end{enumerate}

\subsection{Instalar dependencias}

Una vez que tenemos el proyecto ya añadido a WAMP nos faltaría añadir todas las dependencias, es decir todas las librerías de PHP que se van a instalar con un solo comando gracias al fichero \textit{composer.json}. Para esto seguiremos los siguientes pasos:

\begin{enumerate}
    \item Abrimos el terminal de windows.
    \item Accedemos al directorio "\textit{cd C:\textbackslash{}wamp64\textbackslash{}www\textbackslash{}TFG\_HUBU\textbackslash{}Laravel}"
    \item Ejecutamos el comando \textit{composer install}.
\end{enumerate}

\subsection{Poner a punto las variables de entorno}

Para el correcto funcionamiento de la aplicación vamos a definir las variables de entorno, las cuales se encargan, entre otras cosas, de establecer la conexión con la base de datos. Para esto seguiremos los siguientes pasos:

\begin{enumerate}
    \item Accedemos al directorio "\textit{cd C:\textbackslash{}wamp64\textbackslash{}www\textbackslash{}TFG\_HUBU\textbackslash{}Laravel}"
    \item Creamos el fichero donde se definen las variables de entorno con el comando "\textit{copy .env.example .env}.
    \item Abrimos el fichero .env con cualquier editor de texto.
    \item Cambiamos las variables de entorno:
    \begin{itemize}
        \item APP\_NAME=TFG\_HUBU
        \item DB\_DATABASE=hubu
    \end{itemize}
    \item Generamos la clave automaticamente con el comando "\textit{php artisan key:generate}"
\end{enumerate}

\subsection{Crear la base de datos}

Lo primero que tendremos que hacer será ejecutar wampserver para tener disponibles todas las herramientas. Una vez ejecutado, vamos a crear una base de datos con el nombre \textbf{hubu} como hemos definido anteriormente en las variables de entorno. Esto lo podemos realizar de dos maneras:

\begin{enumerate}
    \item Mediante la consola MySQL:
    \begin{enumerate}
        \item Para acceder a la consola MySQL seleccionamos el icono de wamp con el clic izquierdo, en el menú seleccionamos MySQL y por último seleccionamos MySQL console como podemos ver en la imagen \ref{fig:mysql_consola}.
        \imagenTam{mysql_consola}{Acceder a la consola MySQL}{0.8}
        \item Una vez en la consola ingresamos el usuario \textbf{root} y la contraseña la dejamos vacía.
        \item Con la sesión ya iniciada ejecutamos el comando "\textit{CREATE DATABASE memoria;}".
    \end{enumerate}
    \item Mediante PhpMyAdmin:
    \begin{enumerate}
        \item Accedemos en el navegador a la url "\textit{http://localhost/phpmyadmin/}".
        \item Iniciamos sesión con el usuario \textbf{root} y la contraseña vacía.
        \item Con la sesión ya iniciada seleccionamos \textbf{Nueva} en el menú (\ref{fig:phpmyadmin_crear}) y le ponemos como nombre a la base de datos \textbf{hubu}.
        \imagenTam{phpmyadmin_crear}{Acceder a la consola MySQL}{0.3}
    \end{enumerate}
\end{enumerate}

\subsection{Poner a punto la base de datos}

Una vez creada la base de datos necesitamos crear la estructura y añadir los datos necesarios para que funcione como los roles y el primer usuario para poder realizar el login. Esto lo podemos realizar de dos maneras:

\begin{enumerate}
    \item Mediante \textit{migrations} y \textit{seeders} de Laravel:
    \begin{enumerate}
        \item Abrimos la terminal de windows, accedemos a la ruta "\textit{C:\textbackslash{}wamp\\64\textbackslash{}www\textbackslash{}TFG\_HUBU\textbackslash{}Laravel}"
        \item Ejecutamos el comando "\textit{php artisan migrate}" para ejecutar las migraciones y crear la estructura de tablas.
        \item Ejecutamos el comando "\textit{php artisan db:seed}" para ejecutar los \textit{seeders} y añadir los primeros datos.
    \end{enumerate}
    \item Mediante scripts SQL:
    \begin{enumerate}
        \item Abrimos la consola de MySQL como hemos explicado en el punto anterior.
        \item Una vez en la consola ingresamos el usuario \textbf{root} y la contraseña la dejamos vacía.
        \item Ejecutamos el comando "\textit{use hubu}" para que los scripts se ejecuten sobre esta base de datos.
        \item Ejecutamos los 3 scripts con los siguientes comandos:
        \begin{itemize}
            \item "\textit{source C:/wamp64/www/TFG\_HUBU/SQL/create\_users\_roles.sql;}"
            \item "\textit{source C:/wamp64/www/TFG\_HUBU/SQL/create\_pacientes.sql;}"
            \item "\textit{source C:/wamp64/www/TFG\_HUBU/SQL/add\_roles\_admin.sql;}"
        \end{itemize}
    \end{enumerate}
\end{enumerate}

\subsection{Acceso a la web}

Una vez se han realizado todos estos pasos ya podemos acceder a la web. Para esto seguiremos los siguientes pasos:

\begin{enumerate}
    \item Accedemos a la url "\textit{http://localhost/TFG\_HUBU/Laravel/public}"
    \item Iniciamos sesión con el usuario creado cuyo correo es \textbf{administrador@gmail.com} y contraseña \textbf{1234}.
\end{enumerate}

\section{Pruebas del sistema}

Como ya se ha comentado en la memoria para este proyecto se han realizado tres tipos de pruebas: pruebas de integración, pruebas de validación y pruebas de seguridad. Como las pruebas de validación son pruebas manuales realizadas tanto por el cliente como por el usuario final no va a existir ninguna manera de ejecutarlas automáticamente como las otras dos.

\subsection{Pruebas de integración}

Las pruebas de integración se han realizado para asegurar que la introducción, la modificación y la eliminación tanto de pacientes como de usuarios funciona sin ningún problema.

Para ejecutar estas pruebas primer se han tenido que realizar todos los pasos indicados en el apartado \ref{manual}, luego se tendrán que seguir los siguientes pasos:

\begin{enumerate}
    \item Abrir el terminal de windows.
    \item Acceder a la ruta "\textit{C:\textbackslash{}wamp64\textbackslash{}www\textbackslash{}TFG\_HUBU\textbackslash{}Laravel}
    \item Ejecutar el comando "\textit{php artisan test --testsuite=Feature}"
\end{enumerate}

\newpage

Si todo se ha realizado correctamente el resultado de los test será el de la imagen \ref{fig:test_integracion}.

\imagenTam{test_integracion}{Resultado de los test de integración}{0.7}

\subsection{Pruebas de seguridad}

Estas pruebas se han realizado para asegurar el correcto funcionamiento de los \textit{middleware} y asegurar que usuarios sin permiso no puedan acceder a rutas restringidas.

Para la ejecución de las pruebas se han de seguir los siguientes pasos:

\begin{enumerate}
    \item Abrir el terminal de windows.
    \item Acceder a la ruta "\textit{C:\textbackslash{}wamp64\textbackslash{}www\textbackslash{}TFG\_HUBU\textbackslash{}Laravel}
    \item Ejecutar el comando "\textit{php artisan test --testsuite=Seguridad}"
\end{enumerate}

Si todo se ha realizado correctamente el resultado de los test será el de la imagen.

\textcolor{red}{Faltan por implementar}
\apendice{Documentación de usuario}

\section{Introducción}

En este anexo se va explicar los requisitos de los usuarios para poder acceder a la aplicación web, la instalación y un manual donde se explicarán todas las funcionalidades.

\section{Requisitos de usuarios}

Como es una aplicación web y esta subida a un servidor los únicos requisitos son:
\begin{itemize}
    \item Tener un navegador actualizado que soporte HTML 5.
    \item Tener una cuenta para poder acceder a través del login.
\end{itemize}

\section{Instalación}

La aplicación no requiere de ningún tipo de instalación ya que esta subida a internet y disponible para todo el mundo accediendo a la url \url{http://tfg-hubu.herokuapp.com/}, está es la versión de copia realizada para el tribunal y para la presentación, la web original no se puede mostrar por asuntos de protección de datos. En caso de querer instalar la aplicación web en local se deberán seguir los pasos del apartado \ref{manualProg}.

\section{Manual del usuario}

En este apartado se va a explicar como utilizar todas las funcionalidades de la aplicación web.

\subsection{Gestión de usuarios}

Para acceder a la gestión de usuarios se ha de iniciar sesión con una cuenta de administrador y se tendrá que seleccionar la sección "Gestionar usuarios"{} en el menú desplegable (\ref{fig:menu_desplegable}).

\imagenTam{menu_desplegable}{Menú desplegable de la web.}{0.2}

\subsubsection{Añadir nuevo usuario}

Una vez en la vista de gestión de usuarios, en el menú de la izquierda seleccionamos el apartado "{}Crear nuevo usuario"{}.

\imagenTam{menu_gestion}{Menú de gestión de usuarios.}{0.4}

Rellenamos todos los campos del usuario que queremos añadir al sistema y clicamos el botón "{}Confirmar".  

\imagenTam{anadir_usuario}{Formulario de creación de usuarios.}{0.9}

En el caso en el que alguno de los datos introducidos no sean correctos, como por ejemplo campos en blanco, correo repetido o contraseñas no coincidentes, el formulario devolverá error en los campos incorrectos.

\subsubsection{Modificar usuario existente}

Una vez en la vista de gestión de usuarios, seleccionamos "{}Editar"{} en el usuario que queremos modificar. 

\imagenTam{gestion_usuario}{Tabla de usuarios.}{1}

Modificamos los datos que queremos cambiar del usuarios y clicamos en botón "{}Confirmar".

\imagenTam{modificar_usuario}{Formulario de modificación de usuarios.}{1}

En este formulario también hay control de errores por si se deja algún campo en blanco o el correo ya existe en la base de datos.

\subsubsection{Eliminar usuario existente}

Para eliminar un usuario primero este usuario ha de tener el rol de Oncólogo, así evitamos que se pueda eliminar sin querer un usuario Administrador, que en el caso que solo exista uno sería un problema bastante grave.

Una vez en la vista de gestión de usuarios, seleccionamos "{}Eliminar"{} en el usuario que queramos eliminar.

\imagenTam{gestion_usuario}{Tabla de usuarios.}{1}

Se va a pedir una confirmación de la eliminación para evitar el error humano en la medida de los posible.

\imagenTam{confirmacion_eliminar}{Confirmación de la eliminación.}{0.8}

\subsection{Gestión de datos personales}

Para acceder a la gestión de datos personales se tendrá que seleccionar la sección "Datos personales"{} en el menú desplegable (\ref{fig:menu_desplegable}).

\imagenTam{menu_desplegable}{Menú desplegable de la web.}{0.2}

\subsubsection{Editar datos personales}

Modificamos los datos que queremos cambiar de nuestra cuenta y clicamos en el botón "{}Confirmar". 

\imagenTam{modificar_datos_personales}{Formulario de modificación de datos personales.}{1}

En este formulario también hay control de errores por si se deja algún campo en blanco o el correo ya existe en la base de datos.

\subsubsection{Editar contraseña}

Una vez en la vista de gestión de datos personales, en el menú de la izquierda seleccionamos el apartado "{}Editar contraseña"

\imagenTam{menu_datos_personales}{Menú de gestión de datos personales.}{0.4}

Rellenamos el formulario con la nueva contraseña y con la contraseña antigua y clicamos en el botón "{}Confirmar" 

\imagenTam{modificar_contrasena}{Formulario de modificación de contraseña.}{1}

Este formulario tiene control de errores para comprobar que la contraseña antigua introducida coincide con la almacenada en la base de datos y que ambas contraseñas nuevas coinciden. 

\subsection{Gestión de pacientes}

Para acceder a la gestión de datos personales se tendrá que seleccionar la sección "Pacientes"{} en el menú desplegable (\ref{fig:menu_desplegable}).

\imagenTam{menu_desplegable}{Menú desplegable de la web.}{0.2}

\subsubsection{Creación de paciente}

Una vez en la vista de gestión de pacientes, en el menú de la izquierda seleccionamos el apartado "{}Añadir nuevo paciente".

\imagenTam{menu_pacientes}{Menú de gestión de usuarios.}{0.4}

Rellenamos el formulario con los datos del paciente que queremos crear.

\imagenTam{anadir_paciente}{Formulario de creación de pacientes.}{1}

Este formulario tiene control de errores para que la fecha de nacimiento no pueda ser posterior a la actual y para no permitir campos en blanco.

\subsubsection{Modificar paciente}

Una vez en la vista de gestión de pacientes, seleccionamos "{}Editar"{} en el paciente que queremos modificar.

Los pacientes que no tienen la opción de editar es porque han fallecido\footnote{Requisito de los oncólogos}.

\imagenTam{listado_pacientes}{Tabla de pacientes.}{1}

Modificamos los datos del paciente que queremos cambiar y clicamos en el botón "{}Guardar". 

\imagenTam{datos_paciente}{Formulario de modificación de paciente.}{1}

En el menú de la izquierda se podrá elegir los atributos a modificar y todos tendrán la misma estructura, cuadros de texto con los datos anteriores por si se desea modificar, botón de eliminar por si se quiere borrar y botón de añadir para añadir datos nuevos.

\imagenTam{modificar_metastasis}{Formulario de modificación de metástasis de paciente.}{1}

\subsubsection{Eliminar paciente}

Una vez en la vista de gestión de pacientes, en el menú de la izquierda seleccionamos el apartado "{}Eliminar paciente".

\imagenTam{menu_pacientes}{Menú de gestión de usuarios.}{0.4}

Seleccionamos "{}Eliminar"{} en el paciente que queramos eliminar.

\imagenTam{eliminar_paciente_vista}{Tabla de pacientes.}{1}

Se va a pedir una confirmación de la eliminación para evitar el error humano en la medida de los posible.

\imagenTam{eliminar_paciente}{Confirmación de la eliminación.}{0.8}

\subsection{Visualización de datos}

\subsubsection{Visualización de datos generales}

Para acceder a la visualización de datos generales se tendrá que seleccionar la sección "Gráficas"{} en el menú desplegable (\ref{fig:menu_desplegable}).

\imagenTam{menu_desplegable}{Menú desplegable de la web.}{0.2}

Una vez en la vista de las gráficas, podremos seleccionar en el primer desplegable el tipo de gráfica y podremos elegir dos divisiones para realizar estas gráficas.

Dando clic a "Nueva división{}"{}, aparecerá un desplegable dando las opciones por las que dividir la gráfica.

\imagenTam{graficas_vista}{Selección de divisiones en las gráficas.}{1}

En este caso hemos seleccionado como divisiones el sexo y el tipo de muestra y como vemos en la imagen \ref{fig:ejemplo_graficas} nos ha devuelto la gráfica y la tabla de frecuencias.

Debajo tanto de la tabla como de la gráfica hay un botón que permite descargarlas y otro botón que permite realizar otra gráfica diferente.

\imagenTam{ejemplo_graficas}{Ejemplo de gráfica de dato general.}{1}

\subsubsection{Visualización de datos numéricos}

Para acceder a la visualización de datos numéricos se tendrá que seleccionar la sección "Graficas"{} en el menú desplegable (\ref{fig:menu_desplegable}).

\imagenTam{menu_desplegable}{Menú desplegable de la web.}{0.2}

Una vez en la vista de gráficas, en el menú de la izquierda seleccionamos el apartado "Percentiles".

\imagenTam{menu_graficas}{Menú lateral gráficas.}{0.4}

Seleccionar el dato que queremos analizar.

\imagenTam{percentiles}{Selección de dato numérico.}{1}

En este caso hemos seleccionado como dato a analizar la edad, y como vemos en la imagen \ref{fig:percentiles_graficas} nos ha devuelto la gráfica y la tabla de percentiles.

Debajo tanto de la tabla como de la gráfica hay un botón que permite descargarlas y otro botón que permite realizar otra gráfica diferente.

\imagenTam{percentiles_graficas}{Ejemplo de gráfica de dato numérico.}{1}

\subsubsection{Visualización de datos gráfica kaplan meier}

Para acceder a la visualización de datos numéricos se tendrá que seleccionar la sección "Graficas"{} en el menú desplegable (\ref{fig:menu_desplegable}).

\imagenTam{menu_desplegable}{Menú desplegable de la web.}{0.2}

Una vez en la vista de gráficas, en el menú de la izquierda seleccionamos el apartado "Kaplan meier".

\imagenTam{menu_graficas}{Menú lateral gráficas.}{0.4}

Seleccionar el dato por el que queremos dividir la gráfica.

\imagenTam{kaplan_meier_seleccion}{Selección de dato numérico.}{1}

En este caso hemos seleccionado como dato por el cual separa la gráfica la edad, y como vemos en la imagen \ref{fig:percentiles_graficas} nos ha devuelto la gráfica kaplan meier sin divisiones y la gráfica kaplan meier con la división seleccionada.

Debajo tanto de ambas gráficas hay un botón que permite descargarlas y otro botón que permite realizar otra gráfica diferente.

\imagenTam{kaplan_meier_graficas}{Ejemplo de gráfica de dato numérico.}{1}

\subsubsection{Visualización de datos individuales de cada paciente}

Para acceder a la visualización de datos individuales de cada paciente se tendrá que seleccionar la sección "Pacientes"{} en el menú desplegable (\ref{fig:menu_desplegable}).

\imagenTam{menu_desplegable}{Menú desplegable de la web.}{0.2}

Una vez en la vista de gestión de usuarios, seleccionamos el botón "Ver"{} en el paciente cuyos datos queremos visualizar.

\imagenTam{listado_pacientes}{Tabla de pacientes.}{1}

Podremos ver tablas con todos los datos del paciente seleccionado.

\imagenTam{datos_individuales}{Vista de datos individuales de un paciente.}{1}

Para ver otros datos del paciente como por ejemplo sus síntomas, sus tratamientos o sus seguimientos en el menú de la derecha podemos seleccionar el atributo que queremos visualizar.

\imagenTam{menu_visualizacion}{Menú lateral de la visualización de datos individuales.}{0.4}

\subsubsection{Exportar datos a fichero Excel}

Para acceder a la exportación de la base de datos se tendrá que seleccionar la sección "{}Exportar datos"{} en el menú desplegable (\ref{fig:menu_desplegable}).

\imagenTam{menu_desplegable}{Menú desplegable de la web.}{0.2}

Escribimos en el cuadro de texto el nombre que se quiere que tenga el fichero Excel donde van a estar los datos de todos los pacientes y clicamos el botón "{}Exportar datos".

\imagenTam{exportar_datos}{Vista para exportar los datos.}{1}

\subsection{Realizar filtro}

Se puede realizar un filtro sobre los pacientes en las vistas de las gráficas y en las vista de pacientes. En cualquiera de estas vistas clicamos el botón de aplicar filtro.

\imagenTam{aplicar_filtro}{Botón de filtro.}{1}

Se abrirá un formulario en cual podremos elegir tanto los biomarcadores como los tratamientos por los cuales queremos realizar el filtro de los pacientes.

\imagenTam{filtrar_datos_vista}{Formulario de filtrado de datos.}{1}

En el caso de la tabla de pacientes, la tabla aparecerá con los pacientes que cumplan ese filtro y en el caso de las gráficas estas se realizarán solo teniendo en cuenta los datos de los pacientes que cumplan ese filtro.

\imagenTam{filtrado_tabla}{Tabla de pacientes filtrados por NGS y quimioterapia.}{1}

\imagenTam{filtrado_grafica}{Gráfica de pacientes filtrados por NGS y quimioterapia.}{1}

\subsection{Crear base de datos sintética}

Para crear los datos sintéticos para nuestra base de datos se tendrá que seleccionar la sección "Base sintética"{} en el menú desplegable (\ref{fig:menu_desplegable}).

\imagenTam{menu_desplegable}{Menú desplegable de la web.}{0.2}

Rellenamos el formulario con la información con la cual queramos crear los datos sintéticos y clicamos el botón "{}Crear base de datos".

Al lado de cada cuadro de texto tenemos una ayuda en la cual si ponemos el ratón encima nos muestra información del rango de valores permitidos.

\imagenTam{base_sintetica}{Formulario para crear los datos sintéticos.}{1}

Este formulario tiene control de errores para no permitir campos en blanco y no permitir valores no validos, como por ejemplo, un valor de p que no este entre el 0 y el 1.


\bibliographystyle{plain}
\bibliography{bibliografiaAnexos}

\end{document}
